% ═══════════════════════════════════════════════════════════════
% ML-01d: ¿Qué tal? (MUSTERLÖSUNG)
% Spanisch Klasse 7 - Sequenz 1: Empezamos
% ═══════════════════════════════════════════════════════════════

\documentclass[11pt, a4paper]{article}

\newif\ifswmodus
\swmodusfalse

\usepackage[utf8]{inputenc}
\usepackage[T1]{fontenc}
\usepackage[ngerman]{babel}
\usepackage{helvet}
\renewcommand{\familydefault}{\sfdefault}

\usepackage[margin=2cm]{geometry}
\usepackage{enumitem}
\usepackage{tabularx}
\usepackage{booktabs}

\usepackage{xcolor}
\usepackage{tcolorbox}
\tcbuselibrary{skins}

\usepackage{fancyhdr}
\usepackage{lastpage}
\usepackage{needspace}

% FARBEN
\definecolor{spanisch-color}{HTML}{c41e3a}
\definecolor{grau-color}{HTML}{888888}
\definecolor{hellgrau-color}{HTML}{f5f5f5}
\definecolor{orange-color}{HTML}{b35c00}
\definecolor{loesung-color}{HTML}{c62828}

\definecolor{sw-dark}{HTML}{000000}
\definecolor{sw-gray}{HTML}{666666}
\definecolor{sw-white}{HTML}{ffffff}

\ifswmodus
    \colorlet{spanisch}{sw-dark}
    \colorlet{grau}{sw-gray}
    \colorlet{hellgrau}{sw-white}
    \colorlet{orange}{sw-dark}
    \colorlet{loesung}{sw-dark}
\else
    \colorlet{spanisch}{spanisch-color}
    \colorlet{grau}{grau-color}
    \colorlet{hellgrau}{hellgrau-color}
    \colorlet{orange}{orange-color}
    \colorlet{loesung}{loesung-color}
\fi

\newcommand{\fachfarbe}{spanisch}

\newcommand{\thema}{¿Qué tal?}
\newcommand{\sequenz}{01}
\newcommand{\block}{d}
\newcommand{\fach}{Spanisch}
\newcommand{\klasse}{7}
\newcommand{\maxpunkte}{20}
\newcommand{\datum}{\today}

\pagestyle{fancy}
\fancyhf{}
\renewcommand{\headrulewidth}{0.4pt}
\renewcommand{\footrulewidth}{0.4pt}

\lhead{\fach{} -- Klasse \klasse}
\chead{\textcolor{loesung}{\textbf{ML-\sequenz\block{} (MUSTERLÖSUNG)}}}
\rhead{\datum}

\lfoot{}
\cfoot{}
\rfoot{Seite \thepage\ von \pageref{LastPage}}

\newcommand{\punktebox}[1]{\hfill\fbox{\small \_/#1}}
\newcommand{\luecke}[1]{\rule{#1}{0.4pt}}
\newcommand{\lsg}[1]{\textcolor{loesung}{\textbf{#1}}}

\newtcolorbox{merkkasten}[1][]{
    colback=hellgrau,
    colframe=\fachfarbe,
    colbacktitle=white,
    coltitle=\fachfarbe,
    boxrule=1pt,
    arc=0pt,
    left=8pt, right=8pt, top=8pt, bottom=8pt,
    fonttitle=\bfseries,
    attach boxed title to top left={yshift=-2mm, xshift=5mm},
    boxed title style={boxrule=0pt, colback=white},
    title=#1
}

\newtcolorbox{hinweis}{
    colback=white,
    colframe=orange,
    colbacktitle=white,
    coltitle=orange,
    boxrule=1pt,
    arc=0pt,
    left=5pt, right=5pt, top=5pt, bottom=5pt,
    fonttitle=\bfseries,
    attach boxed title to top left={yshift=-2mm, xshift=5mm},
    boxed title style={boxrule=0pt, colback=white},
    title=Hinweis
}

\newenvironment{aufgabe}[1]{%
    \needspace{5\baselineskip}%
    \nopagebreak[4]%
    \vspace{10pt}
    \noindent\textbf{\textcolor{\fachfarbe}{Aufgabe #1}}
    \nopagebreak[4]%
    \vspace{5pt}
    \nopagebreak[4]%
    \begin{list}{}{\leftmargin=0pt}
    \item[]
}{%
    \end{list}
}

\newcommand{\es}[1]{\textcolor{\fachfarbe}{\textbf{#1}}}

\begin{document}

\begin{center}
    {\Large\bfseries\textcolor{\fachfarbe}{Arbeitsblatt: \thema}}\\[0.3em]
    {\large\textcolor{loesung}{\textbf{--- MUSTERLÖSUNG ---}}}
\end{center}

\vspace{5pt}

\noindent
\begin{tabularx}{\textwidth}{|X|X|X|}
\hline
\textbf{Name:} \lsg{Musterschüler/in} &
\textbf{Klasse:} \klasse &
\textbf{Datum:} \datum \\
\hline
\end{tabularx}

\vspace{15pt}

\begin{merkkasten}[Merkkasten: Nach dem Befinden fragen]

\begin{tabular}{ll}
\textbf{Fragen:} & \\[0.3em]
\es{¿Qué tal?} & = \lsg{Wie geht's?} \\[0.5em]
\es{¿Cómo estás?} & = \lsg{Wie geht es dir?} \\[0.8em]
\textbf{Antworten:} & \\[0.3em]
\es{Muy bien, gracias.} & = \lsg{Sehr gut, danke.} (++) \\[0.5em]
\es{Bien, gracias.} & = \lsg{Gut, danke.} (+) \\[0.5em]
\es{Regular. / Así así.} & = \lsg{Es geht so.} (0) \\[0.5em]
\es{Mal.} & = \lsg{Schlecht.} (--) \\[0.8em]
\textbf{Rückfrage:} & \\[0.3em]
\es{¿Y tú?} & = \lsg{Und du?} \\
\end{tabular}

\end{merkkasten}

\vspace{10pt}

\begin{aufgabe}{1 -- Ordne zu: Wie geht es der Person?}
Verbinde die Antwort mit der passenden Stimmung. \punktebox{4}
\end{aufgabe}

\begin{center}
\begin{tabular}{rcl}
\es{Muy bien, gracias.} & \lsg{---} & sehr schlecht (--) \lsg{$\rightarrow$ Muy bien = ++} \\[0.8em]
\es{Regular.} & \lsg{---} & sehr gut (++) \lsg{$\rightarrow$ Regular = 0} \\[0.8em]
\es{Mal.} & \lsg{---} & es geht so (0) \lsg{$\rightarrow$ Mal = --} \\[0.8em]
\es{Bien, gracias.} & \lsg{---} & gut (+) \lsg{$\rightarrow$ Bien = +} \\
\end{tabular}
\end{center}

\vspace{10pt}

\begin{aufgabe}{2 -- Vervollständige den Dialog}
A und B treffen sich. \punktebox{6}
\end{aufgabe}

\noindent
\textbf{A:} ¡Hola! \lsg{¿Qué tal?}

\vspace{0.8em}

\noindent
\textbf{B:} Bien, gracias. \lsg{¿Y tú?}

\vspace{0.8em}

\noindent
\textbf{A:} Muy bien. ¿Cómo te llamas?

\vspace{0.8em}

\noindent
\textbf{B:} \lsg{Me llamo [Name]. / Soy [Name].}

\vspace{0.8em}

\noindent
\textbf{A:} ¡Mucho gusto! Adiós.

\vspace{0.8em}

\noindent
\textbf{B:} \lsg{¡Hasta luego! / ¡Adiós! / ¡Igualmente!}

\vspace{15pt}

\begin{aufgabe}{3 -- Wiederholung: Das Verb \es{ser}}
Ergänze die Konjugation. \punktebox{4}
\end{aufgabe}

\begin{center}
\begin{tabular}{|l|l|}
\hline
\textbf{Person} & \textbf{ser} \\
\hline
yo & \lsg{soy} \\
\hline
tú & \lsg{eres} \\
\hline
él / ella & \lsg{es} \\
\hline
nosotros & \lsg{somos} \\
\hline
\end{tabular}
\end{center}

\vspace{10pt}

\begin{aufgabe}{4 -- Wiederholung: Artikel einsetzen}
Wähle: \es{el}, \es{la}, \es{los} oder \es{las}. \punktebox{6}
\end{aufgabe}

\noindent
a) \lsg{el} libro \hfill
b) \lsg{la} mesa \\[0.8em]
c) \lsg{los} chicos \hfill
d) \lsg{la} profesora \\[0.8em]
e) \lsg{las} sillas \hfill
f) \lsg{el} cuaderno \\

\vspace{20pt}
\hrule
\vspace{10pt}

\begin{flushright}
\textbf{Punkte:} \lsg{\maxpunkte} / \maxpunkte
\end{flushright}

\end{document}
